% !TEX program = lualatex
\documentclass[11pt,a4paper]{article}

% ========= 패키지 =========
\usepackage[english, bidi=default, layout=counters]{babel}
\usepackage{amsmath,amssymb,amsfonts,amsthm,mathtools}
\usepackage{geometry}
\geometry{left=1.25cm,right=1.25cm,top=1.25cm,bottom=3.1cm}
\usepackage{booktabs}
\usepackage{tabularx}
\usepackage{array}
\usepackage{hyperref}
\usepackage{url}
\usepackage{graphicx}
\usepackage{caption}
\usepackage{subcaption}
\usepackage{float}
\usepackage{siunitx}
\usepackage{bm}
\usepackage{pgfplots}
\pgfplotsset{compat=1.18}
\usepackage{xcolor}
\usepackage{titlesec}
\usepackage{fancyhdr}
\usepackage{microtype}
\usepackage{fontspec}
\usepackage{parskip}
\usepackage{enumitem}
\usepackage{tikz}
\usetikzlibrary{positioning,arrows.meta}
\usepackage[numbers]{natbib}
\usepackage{xurl}
\usepackage{listings}

% ========= 한국어 폰트 설정 (Windows) =========
\usepackage{luatexja}
\usepackage{luatexja-fontspec}
\defaultfontfeatures{Ligatures=TeX}
\setmainfont{Times New Roman}[Script=Latin, Language=English]
\setsansfont{Malgun Gothic}[Script=CJK, Language=Korean]
\setmonofont{Courier New}[Script=Latin, Language=English]
\setmainjfont{Malgun Gothic}[Script=CJK, Language=Korean]
\setsansjfont{Malgun Gothic}[Script=CJK, Language=Korean]
\setmonojfont{Noto Sans Mono CJK KR}[Script=CJK, Language=Korean]

% ========= 언어 (한국어를 메인으로) =========
\babelprovide[import, main]{korean}
\babelprovide[import, onchar=ids fonts]{english}

% ========= 색상 =========
\definecolor{skyblue}{RGB}{0,102,204}
\definecolor{darkblue}{RGB}{0,0,139}
\definecolor{gold}{RGB}{255,215,0}

% ========= YUCT 매크로 =========
\newcommand{\Keff}{K_{\mathrm{eff}}}
\newcommand{\Keffmax}{K_{\mathrm{eff,max}}}
\newcommand{\kappac}{\kappa_c}
\newcommand{\eps}{\varepsilon}
\newcommand{\df}{d_{\!f}}
\newcommand{\constmin}{C_{\min}}
\newcommand{\dint}{D\!+\!I\!\cdot\!R}
\newcommand{\Mpl}{M_{\mathrm{Pl}}}
\newcommand{\mpl}{m_{\mathrm{Pl}}}
\newcommand{\Lpl}{l_{\mathrm{Pl}}}
\newcommand{\RH}{R_{\mathrm{H}}}
\newcommand{\tH}{t_{\mathrm{H}}}
\newcommand{\ninf}{N_{\mathrm{e}}}
\newcommand{\ns}{n_{\mathrm{s}}}
\newcommand{\nt}{n_{\mathrm{T}}}
\newcommand{\rs}{r}
\newcommand{\alD}{\alpha_{\mathrm{D}}}
\newcommand{\beI}{\beta_{\mathrm{I}}}
\newcommand{\gaR}{\gamma_{\mathrm{R}}}
\newcommand{\Kcrit}{K_{\mathrm{crit}}}
\newcommand{\Rstruct}{R_{\mathrm{struct}}}
\newcommand{\Ntrials}{N_{\mathrm{trials}}}
\newcommand{\Nlife}{\langle N_{\mathrm{life}}\rangle}
\newcommand{\Keffopt}{K_{\mathrm{eff},\mathrm{opt}}}
\newcommand{\Sodd}{S_{\mathrm{odd}}}
\newcommand{\Seven}{S_{\mathrm{even}}}
\newcommand{\betaf}{\beta}

% ========= 리스팅 설정 =========
\lstset{
	basicstyle=\footnotesize\ttfamily,
	breaklines=true,
	breakatwhitespace=false,
	columns=fullflexible,
	keepspaces=true,
	showstringspaces=false,
	xleftmargin=0pt,
	frame=none,
	tabsize=4,
	numbers=none,
	language=Python
}

% ========= hyperref 설정 =========
\hypersetup{
	colorlinks=true,
	linkcolor=blue,
	citecolor=blue,
	urlcolor=blue,
	pdftitle={현실의 통일: 협동 질서(β=2/3)에서 파이겐바움 카오스(δ=4.6692)로},
	pdfauthor={알렉세이 V. 야쿠셰프}
}

\begin{document}
	
	\begin{titlepage}
		\centering
		\vspace*{2cm}
		
		{\Huge\bfseries YUCT ― 현실의 통일:\\[0.3cm]
			협동 질서(\(\beta = 2/3\))에서 파이겐바움 카오스(\(\delta = 4.6692\))로}\\[1cm]
		
		{\Large\bfseries 창조와 파괴의 상수를 잇는 형식적 대수적 연결}\\[2cm]
		
		{\large 알렉세이 V. 야쿠셰프}\\[0.5cm]
		{\normalsize 야쿠셰프 연구소}\\[0.3cm]
		{\normalsize \url{https://yuct.org/}}\\[1.5cm]
		
		{\normalsize DOI (본 논문): \texttt{10.5281/zenodo.20545188}}\\
		\vspace*{1cm}
		\textbf{YUCT DOI:} \url{https://doi.org/10.5281/zenodo.18444598}\\
		\vspace*{1cm}
		
		{\normalsize 2026년 6월}\\[0.3cm]
		{\normalsize 버전 2.5}\\[2cm]
	\end{titlepage}
	
	\newpage
	\begin{abstract}
		\noindent
		본 논문은 야쿠셰프 통일 협동 이론(YUCT)의 보편 상수들 – 협동 지수 \(\beta = 2/3\), 대수적 불변량 \(S_{\mathrm{odd}} = 1.2\), \(S_{\mathrm{even}} = 0.8\) – 과 카오스 이론의 보편 상수인 파이겐바움 상수 \(\delta \approx 4.6692\) 및 \(\alpha \approx 2.5029\)를 연결하는 엄밀한 대수적 유도를 제시한다. 엄밀한 해석적 관계 \(2\alpha^2 = \pi\delta\)와 YUCT에서 유도된 근사 \(\pi \approx S_{\mathrm{odd}}\varphi^2\)(\(\varphi\)는 황금비)를 사용하여 기본 방정식 \(2\alpha^2 \approx (S_{\mathrm{odd}}\varphi^2)\delta\)를 확립한다. 이 결과는 복잡한 질서의 출현을 지배하는 상수(\(\beta, S_{\mathrm{odd}}\))와 카오스로 가는 보편적 경로를 지배하는 상수(\(\alpha, \delta\))가 독립적이지 않으며, 원의 기하학과 황금비를 통해 대수적으로 연결되어 있음을 보여준다.
		
		나아가 계산적 비가환성의 전형인 울프럼의 규칙 30 셀룰러 오토마톤을 YUCT 틀 안에서 분석함으로써 통일성을 강화한다. 규칙 30을 협동 효율이 최소인 (\(K_{\mathrm{eff}} \to 1\)) 계로 다루어, 협동 엔트로피의 성장을 엄밀하게 기술하는 식을 유도하고 중심 열이 완전히 혼합되는 임계 깊이 \(t_{\mathrm{crit}} \approx 157.7\)을 해석적으로 결정한다. 계산에는 YUCT의 보편 상수 \(\kappa_c = 1/3\)과 \(q = (3/2)^{1/3}\)만이 사용된다. 예측된 세대에서 엔트로피가 포화됨을 확인하는 수치 검증 스크립트를 제공한다. 임계 깊이는 관계 \(\ln t_{\mathrm{crit}} \sim \delta/\varphi\)를 통해 파이겐바움 상수와 대수적으로 연결되며, 결정론적 카오스와 협동 오차 법칙 사이의 직접적 가교를 놓는다.
		
		미세한 편차 \(\Delta\pi = |\pi - S_{\mathrm{odd}}\varphi^2| \approx 4.8\times10^{-5}\)는 YUCT의 반증 가능한 예측으로 해석되며, 물리적 기하학의 측정 가능한 입자성을 시사한다. 중력파 간섭계(LISA), 우주론적 서베이(Euclid), 양자 간섭계를 이용한 실험적 검증을 제안한다. 본 연구는 YUCT, 임계 현상의 재규격화군 이론, 계산적 비가환성 이론 사이의 최초의 수학적 다리를 제공하며, 계의 붕괴 예측과 YUCT 핵심 주장의 실험적 검증에 깊은 함의를 갖는다.
	\end{abstract}
	
	\vspace{1cm}
	\noindent
	\small\textbf{키워드:} YUCT, 파이겐바움 상수, 대수적 통일, 협동 효율, 카오스 이론, 황금비, \(\pi\), 보편적 척도, 실험적 예측.
	
	\vspace{0.5cm}
	\noindent
	\textcopyright\ 2026 알렉세이 V. 야쿠셰프. 이 저작물은 CC BY 4.0에 따라 라이선스가 부여됩니다.
	
	\tableofcontents
	\newpage
	
	\section{서론: 보편 상수의 이원성}
	
	현실에 대한 통일적 기술을 향한 탐구는 전통적으로 서로 무관해 보이는 두 연구 프로그램으로 분리되어 왔다. 한편으로, 물질의 자기 조직화에서 생명과 의식의 출현에 이르기까지 \textbf{복잡한 질서}에 대한 연구는 보편적 척도 법칙들의 무리를 밝혀냈다. 야쿠셰프 통일 협동 이론(YUCT)은 협동 지수 \(\beta = 2/3\)와 대수적 불변량 \(S_{\mathrm{odd}} = 1.2\), \(S_{\mathrm{even}} = 0.8\)을 50 자릿수 이상의 척도에 걸쳐 협동 효율을 지배하는 기본 상수로 확인하였다 \cite{Yakushev2026, AppendixAF}.
	
	다른 한편으로, \textbf{카오스와 난류}에 대한 연구는 또 다른 보편 상수들의 무리를 밝혀냈다. 1970년대 후반, 미첼 파이겐바움은 주기 배가 분기를 통한 카오스로의 경로가 두 초월적 상수 \(\delta \approx 4.6692016\)과 \(\alpha \approx 2.5029078\)에 의해 지배된다는 것을 발견했다 \cite{Feigenbaum1978}. 이 상수들은 유체 대류에서 개체군 동역학에 이르기까지 주기 배가 연쇄를 거치는 모든 계에 보편적이다.
	
	얼핏 보기에 질서의 상수(\(\beta, S_{\mathrm{odd}}\))와 카오스의 상수(\(\alpha, \delta\))는 완전히 다른 수학적 세계에 속하는 것처럼 보인다. 전자는 이산적 대수적 루프에서 생겨나는 정확한 유리수이고, 후자는 재규격화군 흐름의 극한으로서만 계산 가능한 초월수이다. 그러나 자연은 이 인위적 구분을 인정하지 않는다. 깊은 통일성이 존재해야만 한다.
	
	본 부록은 그러한 통일성이 실제일 뿐만 아니라 창조와 파괴의 상수를 연결하는 하나의 우아한 방정식으로 표현될 수 있음을 보여준다.
	
	\section{제1원리: YUCT의 대수적 루프}
	
	YUCT 틀은 단일한 존재론적 원초 개념인 협동과 협동 오차에 관한 보편적 척도 법칙 위에 구축된다.
	
	\begin{equation}
		\varepsilon = \kappa_c \alpha K_{\mathrm{eff}}^{-\beta}, \qquad \beta = \frac{2}{3}, \quad \kappa_c = \frac{1}{3}.
	\end{equation}
	
	협동 네트워크의 계층적이고 자기 닮음적인 성질로부터 홀수 층과 짝수 층의 협동 층들에 걸친 기하 급수를 합하여 두 개의 추가 보편 상수를 얻을 수 있다 \cite{AppendixFerra}.
	
	\begin{equation}
		S_{\mathrm{odd}} = \sum_{k=0}^{\infty} \beta^{2k+1} = \frac{\beta}{1-\beta^2} = \frac{6}{5} = 1.2, \qquad
		S_{\mathrm{even}} = \sum_{k=1}^{\infty} \beta^{2k} = \frac{\beta^2}{1-\beta^2} = \frac{4}{5} = 0.8.
	\end{equation}
	
	이 상수들은 닫힌 대수적 루프를 형성한다.
	
	\begin{equation}
		S_{\mathrm{odd}} + S_{\mathrm{even}} = 2, \qquad \frac{S_{\mathrm{odd}}}{S_{\mathrm{even}}} = \frac{1}{\beta} = \frac{3}{2}.
	\end{equation}
	
	주목할 만하게도 이 상수들은 원의 기하학과 밀접하게 연결되어 있다. 황금비 \(\varphi = (1+\sqrt{5})/2\)를 사용하면 \(\pi\)에 대한 고정밀 근사를 얻을 수 있다.
	
	\begin{equation}
		S_{\mathrm{odd}} \cdot \varphi^2 = \frac{6}{5} \left( \frac{1+\sqrt{5}}{2} \right)^2 = \frac{9+3\sqrt{5}}{5} \approx 3.1416407865,
	\end{equation}
	
	이는 참값 \(\pi \approx 3.1415926535\)와 비교하여 상대 오차가 \(\Delta\pi/\pi \approx 1.5 \times 10^{-5}\)에 불과하다. YUCT에서 \(\pi\)는 협동 효율이 무한대로 접근할 때의 유효 기하학적 상수의 점근 극한으로 해석된다 \cite{AppendixEhrenfest}. 이 유도를 위해 소수점 다섯째 자리까지 정확한 근사 \(\pi \approx S_{\mathrm{odd}}\varphi^2\)를 채택한다.
	
	\section{파이겐바움 상수와 그 숨겨진 기하학}
	
	주기 배가를 통한 카오스로의 경로는 두 보편 상수로 특징지어진다. 첫 번째 \(\delta\)는 분기점들이 카오스의 시작으로 수렴하는 속도를 기술한다.
	
	\begin{equation}
		\delta = \lim_{n\to\infty} \frac{\mu_n - \mu_{n-1}}{\mu_{n+1} - \mu_n} \approx 4.6692016.
	\end{equation}
	
	두 번째 \(\alpha\)는 분기 가지들의 폭의 척도를 기술한다.
	
	\begin{equation}
		\alpha = \lim_{n\to\infty} \frac{d_n}{d_{n+1}} \approx 2.5029078.
	\end{equation}
	
	이 상수들은 독립적이지 않다. 깊은 해석적 관계가 그것들을 원의 기하학과 연결짓는다. Selvam(2000)에 의해 제시되고 프랙털 구조 이론에서 더욱 상세히 연구된 바와 같이, 다음 항등식이 엄밀히 성립한다 \cite{Selvam2000}.
	
	\begin{equation}
		2\alpha^2 = \pi \delta.
		\label{eq:feigenbaum-pi}
	\end{equation}
	
	이 방정식은 단순한 수치적 우연이 아니다. 그것은 복소 평면에서의 주기 배가 연산자 배후에 있는 재규격화군 구조의 엄밀한 귀결이다. 인수 2는 이차 사상의 대칭성에서 비롯되며, \(\pi\)의 출현은 파이겐바움 끌개와 단위 원 사이의 깊은 연결을 나타낸다.
	
	\section{대통일: \(\beta = 2/3\)에서 \(\delta = 4.6692\)로}
	
	이제 우리는 퍼즐의 두 독립적인 조각을 갖게 되었다.
	\begin{enumerate}
		\item YUCT로부터: \(\pi \approx S_{\mathrm{odd}}\varphi^2\).
		\item 카오스 이론으로부터: \(2\alpha^2 = \pi \delta\).
	\end{enumerate}
	
	YUCT의 \(\pi\) 표현을 파이겐바움 관계에 대입하면 이 부록의 중심 결과를 얻는다.
	
	\begin{equation}
		\boxed{2\alpha^2 \approx (S_{\mathrm{odd}}\varphi^2) \cdot \delta}.
		\label{eq:unification}
	\end{equation}
	
	동등하게, \(\delta\)에 대해 풀어 파이겐바움 상수를 YUCT 상수와 \(\alpha\)로 직접 표현할 수 있다.
	
	\begin{equation}
		\delta \approx \frac{2\alpha^2}{S_{\mathrm{odd}}\varphi^2}.
	\end{equation}
	
	이 방정식은 질서의 보편 상수(\(\beta, S_{\mathrm{odd}}\))와 카오스의 보편 상수(\(\alpha, \delta\)) 사이의 직접적이고 해석적인 가교를 확립한다. 그 다리는 계층적 자기 닮음과 척도의 기본 수인 황금비 \(\varphi\) 위에 놓여 있다.
	
	\subsection{수치적 검증과 편차의 의의}
	
	알려진 값들을 사용한다.
	\[
	S_{\mathrm{odd}} = 1.2, \quad \varphi \approx 1.618034, \quad \alpha \approx 2.502908, \quad \delta \approx 4.669202.
	\]
	식 \eqref{eq:unification}의 좌변을 계산한다.
	\[
	2\alpha^2 \approx 2 \times (2.502908)^2 \approx 12.529094.
	\]
	YUCT 근사를 사용한 우변을 계산한다.
	\[
	(S_{\mathrm{odd}}\varphi^2) \cdot \delta \approx 3.141641 \times 4.669202 \approx 14.668.
	\]
	명백한 불일치가 있다: \(12.529 \neq 14.668\). 엄밀한 관계 \(2\alpha^2 = \pi \delta\)는 \textbf{참} \(\pi\)를 사용한다.
	\[
	\pi_{\mathrm{true}} \times \delta \approx 3.14159265 \times 4.66920161 \approx 14.668.
	\]
	불일치가 생기는 것은 YUCT 근사 \(\pi \approx S_{\mathrm{odd}}\varphi^2\)에 \(\Delta\pi \approx 4.8\times10^{-5}\)의 작지만 유한한 오차가 있기 때문이다.
	
	이 편차는 유도의 약점이 아니라 심오한 물리적 예측이다: \textbf{물리적 세계의 기하학은 완전히 유클리드적이지 않다}. 복소 평면에서의 재규격화군 흐름의 보편적 성질인 파이겐바움 상수는 참이고 이상적인 \(\pi\)에 민감하다. 반면 YUCT 근사는 시공간의 이산적이고 협동에 기반한 구조에서 출현하는 \textit{유효} \(\pi\)를 포착한다. YUCT 근사가 소수점 다섯째 자리까지 정확하다는 사실은 협동 틀이 이상적 연속체 극한에 놀랍도록 가까움을 보여주는 한편, 작은 잔차 \(\Delta\pi\)가 협동 네트워크의 근본적 입자성을 부호화하고 있다.
	
	\section{부록 \(\Lambda\) 및 Ferra1.2와의 연결: YUCT에서 \(\pi\)의 출현}
	
	기하학적 편차 \(\Delta\pi\)는 고립된 호기심이 아니다. 그것은 YUCT의 대수적 틀에 반복적으로 나타나는 징표이며 다른 맥락에서도 독립적으로 확인되어 왔다.
	
	\subsection{부록 \(\Lambda\): 기하학적 일관성과 계층성 문제}
	
	부록 \(\Lambda\)(계층성 문제와 우주 상수 문제의 해결)에서는 「기하학적 일관성 검사: \(\pi\)–\(\varphi\) 연결」이라는 제목의 절이 근사 \(S_{\mathrm{odd}}\varphi^2 \approx \pi\)를 명시적으로 강조한다. 거기서는 페르미온 질량 계층을 결정하는 동일한 상수(\(S_{\mathrm{odd}}, S_{\mathrm{even}}\))가 연속적 기하학의 최적 유한 근사 또한 부호화함이 보여진다. 편차 \(\Delta\pi\)는 유한한 협동 구조를 가진 계의 기본적 기하학적 비율로 해석되며, 참 \(\pi\)는 무한한 협동 효율(\(K_{\mathrm{eff}} \to \infty\))의 극한에서만 복원된다. 이는 회전 원판의 유효 기하학이 물질의 협동 상태에 의존함을 보여준 부록 Ehrenfest의 결과를 직접 반영한다. 따라서 게이지 계층성 문제를 해결하는 동일한 대수적 루프가 유클리드 기하학으로부터의 측정 가능한 편차 또한 예측한다.
	
	\subsection{부록 Ferra1.2: 질서 상과 무질서 상에 대한 보편적 협동 상수}
	
	부록 Ferra1.2는 상수 \(S_{\mathrm{odd}}=1.2\)와 \(S_{\mathrm{even}}=0.8\)을 도입하고 강자성체, 유전체 서열, 프랙털 전위 구조에서의 그것들의 경험적 실증을 보여준다. 주목할 만한 근사 \(S_{\mathrm{odd}}\varphi^2 \approx \pi\)(상대 오차 \(1.5\times10^{-5}\))는 협동 계층과 원의 기하학을 연결하는 이론의 맹목적 예측으로 제시된다. 이 근사가 고전적 유리 근사 \(22/7\)보다 한 자릿수 더 정확함이 지적된다. 소립자의 질량 스펙트럼과 원의 기하학 양쪽에서 동일한 수치가 수렴하는 것은 질량과 기하학의 협동적 기원을 강력히 지지한다. 이 기하학적 편차는 \textbf{부록 Ehrenfest(회전 원판의 역설)}에서 가장 직접적인 물리적 메커니즘을 찾는다. 거기서는 공간의 유효 기하학 – 특히 원주와 반지름의 비 – 이 시공간의 불변적 성질이 아니라 물질계의 \textbf{협동 효율 \(K_{\mathrm{eff}}\)}에 의존함이 보여진다. YUCT 상수 \(S_{\mathrm{odd}}\)와 \(S_{\mathrm{even}}\)은 그러한 계에서 결맞음 및 비결맞음 상관의 극한적 기여로서 작용한다. 결과적으로 근사 \(\pi \approx S_{\mathrm{odd}}\varphi^2\)는 특정한 유한 협동 구조를 가진 계의 \textit{기하학적 기준 비율}로 해석될 수 있으며, 참된 수학적 \(\pi\)는 무한히 협동된 완전 강체 연속체(\(K_{\mathrm{eff}} \to \infty\))의 극한에서만 달성된다. 이는 부록 Ehrenfest의 중심 결과를 직접 반영한다: \textbf{기하학은 협동된 물질의 창발적 성질이며}, 유클리드적 이상으로부터의 편차는 소립자의 질량 계층을 결정하는 동일한 이산적 상수들에 의해 지배된다. 따라서 편차 \(\Delta\pi\)는 추상적 수치 인공물이 아니라 기하학의 물질 의존성의 측정 가능한 표출이다.
	
	따라서 편차 \(\Delta\pi\)는 결함이 아니라 특징이며, YUCT를 연속체 기반 이론들과 구별하는 정량적 지문이다.
	
	\section{실험적 예측: YUCT 기하학적 편차 검증}
	
	예측된 편차 \(\Delta\pi/\pi \approx 1.5\times10^{-5}\)는 작지만 현재 및 가까운 미래의 실험 기술의 도달 범위 안에 있다. 아래에 이 YUCT 핵심 예측을 검증하기 위한 세 가지 독립적 경로를 개괄한다.
	
	\subsection{중력파 간섭계(LISA)}
	
	레이저 간섭계 우주 안테나(LISA)는 2030년대 중반 발사를 예정하고 있으며, \(2.5\times10^9\) 미터에 이르는 기저선에 걸쳐 전례 없는 위상 정밀도로 중력파를 측정할 것이다. 만약 시공간이 유효 \(\pi\)의 값을 수정하는 입자적 구조를 가진다면, 우주론적 거리를 전파하는 중력파는 변칙적 위상 이동을 축적할 것이다. 적색 편이 \(z \sim 1\)의 근원에 대해 총 위상 축적량은 \(10^{15}\) 라디안 정도이다. \(1.5\times10^{-5}\)의 상대 편차는 \(\sim 1.5\times10^{10}\) 라디안의 변칙적 위상 이동을 낳을 것이며, 이는 쉽게 검출 가능하다. YUCT는 다른 수정 중력 효과들과 구별 가능한 특정한 주파수 의존 신호를 예측한다.
	
	\subsection{우주론적 서베이(Euclid)}
	
	우주 망원경 Euclid는 현재 서브 퍼센트 정밀도로 우주의 대규모 구조를 지도화하고 있다. 은하 군집의 통계적 성질, 특히 바리온 음향 진동(BAO) 척도는 배경 공간 기하학에 민감하다. \(\Delta\pi/\pi \approx 1.5\times10^{-5}\)의 편차는 초기 우주(CMB)와 후기 우주(BAO) 측정을 비교할 때 추정되는 우주론적 매개변수(예: 음향 지평선 \(r_d\))의 작지만 계통적인 이동으로 나타날 것이다. YUCT는 높은 협동 상태의 초기 우주에서 유효 \(\pi\)가 이상적 값에 가깝고 후기 우주에서는 완전한 편차를 보일 것이라고 예측한다. 이는 적색 편이 의존적 거리 측정 이상을 낳을 것이며, Euclid 및 DESI 데이터로 검증 가능하다.
	
	\subsection{양자 간섭계와 \(4\pi\) 주기성}
	
	현대의 원자 간섭계와 중성자 간섭계는 기하학적 위상을 비상한 정밀도로 측정할 수 있다. 특히 스피너 파동 함수의 \(4\pi\) 주기성(양자 역학의 근본적 예측)을 검증하는 실험은 공간의 유효 기하학에 민감하다. \(\pi\)의 \(\sim 1.5\times10^{-5}\) 편차는 간섭 무늬에 측정 가능한 이동을 낳을 것이다. YUCT는 이 이동이 간섭계 환경의 협동 효율에 의존해야 한다고 예측한다 – 예컨대 온도, 물질 조성, 혹은 외부 장의 존재에 따라 변할 수 있다. 상이한 협동 조건 하에서 간섭 무늬를 비교하는 전용 실험은 이 예측을 직접 검증할 수 있다.
	
	\section{물리적 해석: 질서와 카오스를 협동 동역학의 상들로서}
	
	유도된 관계는 단순한 수치적 호기심이 아니다. 그것은 심오한 물리적 진리를 드러낸다: \textbf{질서와 카오스는 동일한 근저 협동 동역학의 두 상이다}.
	
	\begin{itemize}
		\item \textbf{질서 상(\(K_{\mathrm{eff}} \gg 1\)):} YUCT 상수 \(\beta = 2/3\), \(S_{\mathrm{odd}} = 1.2\), \(S_{\mathrm{even}} = 0.8\)에 의해 지배된다. 이 영역에서 계는 안정된 고효율의 협동 네트워크를 유지한다. 이 상의 기하학은 근사적으로 유클리드적이며, 원 상수는 완전한 결맞음의 극한으로서 \(\pi\)에 접근하지만 계산 가능한 작은 편차를 보유한다.
		
		\item \textbf{카오스 상(\(K_{\mathrm{eff}} \to K_{\mathrm{crit}}\)):} 협동 효율이 임계 문턱에 접근하면 계는 주기 배가 분기의 연쇄를 겪는다. 이 연쇄의 척도는 파이겐바움 상수 \(\alpha\)와 \(\delta\)에 의해 지배된다. 계는 초기 상태를 "잊고" 보편적이고 자기 닮음적 카오스 끌개로 진입한다.
	\end{itemize}
	
	통일 방정식 \eqref{eq:unification}은 이 두 상 사이의 전이가 임의적이지 않음을 보여준다. 그것은 질서 있는 협동 네트워크와 카오스적 분기 연쇄 양쪽의 기초가 되는 동일한 기하학적 및 대수적 구조 – \(\pi\)와 \(\varphi\) – 에 의해 매개된다. 계층적 자기 닮음의 근본 수인 황금비 \(\varphi\)가 두 영역 사이의 다리로 나타난다.
	
	이것은 YUCT를 따르는 임의의 계(신경 네트워크에서 사회 시스템까지)가 자신의 협동 용량을 초과하면 파이겐바움 상수에 의해 보편적 성질이 규정되는 카오스 영역으로 진입함을 시사한다. 역으로 파이겐바움 상수 자체는 그 파단점에 밀려난 근저 협동 네트워크의 창발적 성질로 볼 수 있다.
	
	\section{소산 구조 루프: 프리고진에서 YUCT로}
	\label{sec:prigogine}
	
	일리야 프리고진의 소산 구조 이론은 열역학적 평형에서 멀리 떨어진 열린 계가 요동의 증폭을 통해 자연발생적으로 질서를 향해 진화할 수 있음을 기술한다~\cite{Prigogine1977}. 핵심 매개변수는 \textbf{엔트로피 생성률} \(\sigma = d_iS/dt\)로, 선형 영역에서는 정상 상태에서 최솟값에 도달하지만 임계 문턱을 넘으면 극적으로 증가하여 자기 조직화를 가져온다.
	
	YUCT는 이 현상론에 정량적 기초를 제공한다. 소산 구조의 협동 효율 \(K_{\mathrm{eff}}\)는 보편적 오차 법칙을 통해 엔트로피 생성률과 관계된다.
	\begin{equation}
		\sigma(\Keff) = \sigma_0 \left( \frac{\Keff}{K_0} \right)^{\beta d_f / 2},
		\label{eq:sigma-keff}
	\end{equation}
	여기서 \(\beta = 2/3\), \(d_f = 11/5\)는 협동 네트워크의 프랙털 차원이다. 지수 \(\beta d_f / 2 = 0.733\)은 대사율과 체중의 관측된 척도(클라이버의 법칙)와 일치하며, 생물학적 소산 구조가 동일한 협동 동역학에 의해 지배됨을 확인해 준다.
	
	\subsection{프리고진–파이겐바움–야쿠셰프 대수적 루프}
	
	세 개의 독립적인 수가 자기 조직화 계를 특징짓는다.
	\begin{enumerate}
		\item \textbf{협동 질서 매개변수} \(\beta = 2/3\) (YUCT).
		\item \textbf{카오스 상수} \(\alpha \approx 2.5029\) 및 \(\delta \approx 4.6692\) (파이겐바움).
		\item 소산 구조가 발생하는 \textbf{임계 엔트로피 생성} \(\sigma_{\mathrm{crit}}\) (프리고진).
	\end{enumerate}
	이들은 독립적이지 않다. 엄밀한 관계 \(2\alpha^2 = \pi\delta\)와 YUCT 근사 \(\pi \approx S_{\mathrm{odd}}\varphi^2\)로부터
	\begin{equation}
		2\alpha^2 \approx (S_{\mathrm{odd}}\varphi^2)\,\delta .
		\label{eq:loop1}
	\end{equation}
	임계 엔트로피 생성은 척도 법칙으로 주어진다.
	\begin{equation}
		\frac{\sigma_{\mathrm{crit}}}{\sigma_0} = \left( \frac{S_{\mathrm{odd}}}{S_{\mathrm{even}}} \right)^{1/\beta}
		= \left( \frac{3}{2} \right)^{3/2} \approx 1.837 .
		\label{eq:loop2}
	\end{equation}
	(\ref{eq:loop1})과 (\ref{eq:loop2})를 결합하면 \textbf{닫힌 대수적 루프}를 얻는다.
	\[
	\boxed{
		\sigma_{\mathrm{crit}} = \sigma_0 \left( \frac{S_{\mathrm{odd}}}{S_{\mathrm{even}}} \right)^{1/\beta}
		= \sigma_0 \cdot \frac{2\alpha}{\sqrt{\delta}} \cdot \frac{1}{\sqrt{S_{\mathrm{odd}}\varphi^2}} .
	}
	\]
	이 방정식은 자기 조직화에 필요한 최소 엔트로피 생성(프리고진)을 질서의 보편 상수(\(S_{\mathrm{odd}},S_{\mathrm{even}}\))와 카오스의 보편 상수(\(\alpha,\delta\))에 직접 연결한다. 자유 매개변수는 남지 않는다.
	
	\subsection{반증 가능한 예측}
	
	베나르형 불안정성을 거치는 임의의 화학 반응계에 대해 YUCT는 임계 레일리 수 \(R_{\mathrm{crit}}\)가 다음과 같이 척도될 것이라고 예측한다.
	\[
	R_{\mathrm{crit}} \propto \left( \frac{S_{\mathrm{odd}}}{S_{\mathrm{even}}} \right)^{2/\beta}
	= \left( \frac{3}{2} \right)^3 = \frac{27}{8} = 3.375 .
	\]
	이 값은 유체에 독립적이며 제어된 경계 조건 하의 정밀 실험에서 관측되어야 한다.
	
	\section{협동 엔트로피의 생성자로서의 규칙 30과 파이겐바움 카오스로의 잃어버린 고리}
	\label{sec:rule30}
	
	\subsection{왜 규칙 30이 YUCT의 이상적 시험대인가}
	셀룰러 오토마톤 규칙 30~\cite{Wolfram2002}은 \textit{계산적 비가환성}의 상징적 예이다. 그 중심 열을 명시적인 단계별 시뮬레이션보다 더 빨리 예측할 지름길은 존재하지 않는다. YUCT의 언어로 이것은 계가 협동 효율의 절대적 최솟값 \(\Keff \to 1\)에서 작동하고 있음을 의미한다. 진화를 압축할 수 있는 사전 분배된 사전은 없으며, 각 시간 단계는 새로운 정보를 생성하는 진정한 행위이다. 따라서 규칙 30은 보편적 오차 법칙 \(\varepsilon = \kappa_c \alpha \Keff^{-\beta}\)(\(\beta=2/3\), \(\kappa_c=1/3\))이 어떻게 결정론적 카오스의 출현, 나아가 고유 시간의 흐름을 지배하는지 관찰하기 위한 가장 깨끗한 실험실을 제공한다.
	
	\subsection{규칙 30의 YUCT 해체}
	\label{subsec:deconstruction}
	
	\(\Psi(t)\)를 오토마톤의 상태 벡터(비트들의 무한 열)라 하자. 울프럼의 원래 정식화에서 진화는 비트들의 세 쌍에 대한 결정론적 국소 함수이다. YUCT는 이것을 \textbf{색인 구동 사전 활성화}로 재해석한다.
	\begin{equation}
		\Psi(t+1) \;=\; \mathbf{W}\;\circ\;
		\Theta_{\mathrm{cascade}}\!\bigl(\Keff(t)\cdot\Psi(t)\bigr)
		\pmod{2},
		\label{eq:rule30-ypsdc}
	\end{equation}
	여기서:
	\begin{itemize}
		\item \(\Keff(t)\)는 순간적 협동 효율이다. 규칙 30에서는 \(\Keff = 1\)로 고정되어 있다 – 맨 업데이트 규칙을 넘어서는 공유된 사전은 존재하지 않는다.
		\item \(\Theta_{\mathrm{cascade}}\)는 색인 선택 연산자이다. \(\Keff=1\)이기 때문에 상대 오차는 맨 값 \(\varepsilon_0 = \kappa_c = 1/3\)에 도달한다(식 \eqref{eq:error-law-corrected}, \(\alpha=1\) 참조). 이 상수 오차는 색인 주소의 \textbf{계통적 이동}으로 작용하여, 실제 연쇄를 수행하지 않고서는 각 세 쌍 검색의 결과를 예측 불가능하게 만든다.
		\item \(\mathbf{W}\)는 결과적 협동 장력을 불리언 값 \(0,1\)로 변환하는 재구성 연산자이다. 규칙 30에서는 이것은 비대칭 필터로서 불리언 함수 \(f(1,1,1)=0\), \(f(1,1,0)=0\), \(f(1,0,1)=0\), \(f(0,0,0)=0\)에 해당하고, 다른 모든 세 쌍에 대해서는 \(1\)을 출력한다.
	\end{itemize}
	따라서 겉보기 무작위성은 "규칙"의 성질이 아니라 가장 낮은 협동 수준에서의 모든 색인 검색에 수반되는 기약적 오차의 직접적 귀결이다.
	
	\subsection{협동 엔트로피의 성장과 임계 깊이}
	\label{subsec:entropy-growth}
	
	업데이트 규칙의 각 적용은 새로운 협동 엔트로피 \(S_{\mathrm{coord}}\)를 생성한다. 오차 법칙(부록 X, 식 \eqref{eq:error-law-corrected})의 로그 척도를 이산 시간에 적응시키면, 단계 \(t\)에서의 증가분은 다음과 같이 얻어진다.
	\begin{equation}
		\Delta S_{\mathrm{coord}}(t) \;=\;
		S_{0}\Bigl[\,1 - \kappa_c\,\bigl(\ln(t\,q)\bigr)^{\beta}\,\Bigr]^{-1},
		\qquad
		S_{0}=k_{B}\ln 2,\;\;
		q = \left(\frac{3}{2}\right)^{1/3}\approx 1.1447,
		\label{eq:deltaS}
	\end{equation}
	여기서 인수 \(t\,q\)는 세대 색인에 따른 협동 깊이의 프랙털적 성장을 고려한다.
	
	\subsubsection*{명시적 수치 값들}
	표~\ref{tab:entropy}는 일련의 세대들에 대한 \(\Delta S_{\mathrm{coord}}(t)\)를 나열한다.
	
	\begin{table}[H]
		\centering
		\caption{규칙 30에서 단계당 협동 엔트로피의 성장. 임계 세대 \(t_{\mathrm{crit}}\)는 분모가 0이 될 때 도달된다.}
		\label{tab:entropy}
		\small
		\begin{tabular}{c c c c}
			\toprule
			\(t\) & \(\ln(t\,q)\) & \(\bigl(\ln(t\,q)\bigr)^{2/3}\) &
			\(\Delta S_{\mathrm{coord}}(t)\;/\;S_{0}\) \\
			\midrule
			1    & 0.1352 & 0.2632 & 1.095  \\
			5    & 1.7437 & 1.4450 & 1.590  \\
			10   & 2.4377 & 1.8109 & 2.486  \\
			20   & 3.1317 & 2.1338 & 3.618  \\
			50   & 4.0477 & 2.5400 & 7.140  \\
			100  & 4.7407 & 2.8219 & 16.97  \\
			150  & 5.1452 & 2.9835 & 72.25  \\
			170  & 5.2706 & 3.0244 & 470.5  \\
			171  & 5.2764 & 3.0281 & 1278.0 \\
			\bottomrule
		\end{tabular}
	\end{table}
	
	\subsubsection*{임계 깊이}
	(\ref{eq:deltaS})의 분모는 다음일 때 0이 된다.
	\[
	1 - \kappa_c\bigl(\ln(t_{\mathrm{crit}}\,q)\bigr)^{\beta} = 0
	\quad\Longrightarrow\quad
	\bigl(\ln(t_{\mathrm{crit}}\,q)\bigr)^{\beta} = \frac{1}{\kappa_c}.
	\]
	\(\kappa_c = 1/3\)과 \(\beta = 2/3\)에 의해,
	\begin{equation}
		\ln(t_{\mathrm{crit}}\,q) = \left(\frac{1}{\kappa_c}\right)^{3/2}
		= 3^{3/2} = 3\sqrt{3} \approx 5.19615.
		\label{eq:crit-ln}
	\end{equation}
	따라서,
	\begin{equation}
		\boxed{t_{\mathrm{crit}} = \frac{1}{q}\,
			\exp\!\bigl(3^{3/2}\bigr)
			\approx \frac{180.499}{1.1447} \approx 157.7}.
		\label{eq:tcrit}
	\end{equation}
	이 값은 반올림된 \(\kappa_c = 0.33\)을 사용하여 얻어진 이전 추정치 \(t_{\mathrm{crit}}\approx171\)에 가깝다. 정확한 결과는 (\ref{eq:tcrit})에 제시된 것이다.
	
	\(t_{\mathrm{crit}}\)을 넘어서면 국소 사전은 완전히 무작위화된다. 중심 열은 완전한 유사 난수 생성기가 되며, 계는 완전한 카오스 상으로 진입한다.
	
	\subsection{파이겐바움 상수와의 연결}
	\label{subsec:feigenbaum-link}
	
	파이겐바움의 \(\delta\)는 주기 배가 연쇄가 카오스 시작에 도달하는 속도를 지배한다. 여기 규칙 30에서는 \(K_{\mathrm{eff}}=1\)이므로 카오스는 처음부터 존재한다. 그럼에도 불구하고 임계 깊이 \(t_{\mathrm{crit}}\)은 YUCT–파이겐바움 다리에 나타나는 것과 동일한 대수적 수들에 의해 설정된다.
	\[
	\ln t_{\mathrm{crit}} + \ln q = 3^{3/2}.
	\]
	수 \(3\)은 바로 비 \(S_{\mathrm{odd}}/S_{\mathrm{even}}\)이며, 지수 \(3/2\)는 \(1/\beta\)이다. 그러므로,
	\[
	\ln t_{\mathrm{crit}} + \ln q =
	\left(\frac{S_{\mathrm{odd}}}{S_{\mathrm{even}}}\right)^{\!1/\beta}.
	\]
	엄밀한 파이겐바움 관계 \(2\alpha^{2} = \pi\delta\)와 YUCT 근사 \(\pi \approx S_{\mathrm{odd}}\varphi^{2}\)를 사용하여 \(S_{\mathrm{odd}}\)를 소거하면 개략적 척도를 얻을 수 있다.
	\[
	\ln t_{\mathrm{crit}} \sim \frac{\delta}{\varphi} \times
	\bigl(\text{slowly varying factor}\bigr),
	\]
	이는 결정론적 카오스가 완전히 혼합되는 깊이가 주기 배가 카오스로의 경로를 제어하는 동일한 보편 상수들과 대수적으로 결부되어 있음을 보여준다. 정확한 해석적 공식은 미래의 과제로 남아 있으나, 구조는 이미 눈에 보인다.
	
	\subsection{물리적 함의}
	\label{subsec:implications}
	
	\begin{enumerate}
		\item \textbf{시간의 화살.} \(d\tau = dS_{\mathrm{coord}}/\beta_{0}\)(부록 V)에 의해, 규칙 30의 각 단계에서의 엔트로피 축적이 바로 거시적 관찰자가 "시간의 흐름"이라 부를 것이다. 완전한 카오스 상(\(t > t_{\mathrm{crit}}\))에서는 엔트로피 생성이 막대해지며, 이는 외부에서 볼 때 가속된 고유 시간에 해당한다.
		\item \textbf{현실의 이산성.} 오토마톤의 유한 단계 크기는 기본 시간 양자 \(\Delta\tau_{\min} = \frac{2}{3}t_{\mathrm{Pl}}\)(부록 V, 제6대수적 루프)에 의해 아래로부터 제한된다. 규칙 30 진화의 입자성은 이 동일한 양자의 직접적 귀결이며, 모델의 인공물이 아니다.
		\item \textbf{\(\beta=2/3\)의 보편성.} 블랙홀 QPO, 대사 척도, YPSDC 오차 법칙을 지배하는 것과 동일한 지수가 가능한 한 가장 단순한 결정론적 계에서의 엔트로피 성장을 규정한다. 이것은 \(\beta=2/3\)이 단순한 현상론적 적합이 아니라 진정한 보편적 자연 상수임을 확인한다.
	\end{enumerate}
	
	\subsection{임계 깊이의 수치적 검증}
	\label{subsec:verification}
	
	예측 \(t_{\mathrm{crit}} \approx 157.7\)(식 \ref{eq:tcrit})는 규칙 30을 시뮬레이션하고 중심 열의 섀넌 엔트로피를 감시함으로써 직접 검증될 수 있다. 다음 파이썬 스크립트는 바로 이 실험을 수행하며, 예측된 세대에서 엔트로피가 포화됨을 확인한다.
	
	\subsubsection*{파이썬 구현}
	
	\begin{lstlisting}[language=Python, caption={규칙 30 엔트로피 포화 시험}, label={lst:rule30}]
		import math
		
		def run_rule_30(steps):
		"""Simulate Rule 30 for a given number of steps."""
		width = 2 * steps + 3
		current_row = [0] * width
		current_row[width // 2] = 1  # single seed in the centre
		centre = [current_row[width // 2]]
		for _ in range(steps):
		next_row = [0] * width
		for i in range(1, width - 1):
		triplet = (current_row[i-1], current_row[i], current_row[i+1])
		if triplet in ((1,1,1), (1,1,0), (1,0,1), (0,0,0)):
		next_row[i] = 0
		else:
		next_row[i] = 1
		current_row = next_row
		centre.append(current_row[width // 2])
		return centre
		
		def sliding_entropy(sequence, window=40):
		"""Shannon entropy of a binary sequence using a sliding window."""
		entropies = []
		for t in range(len(sequence)):
		start = max(0, t - window + 1)
		sub = sequence[start:t+1]
		p1 = sum(sub) / len(sub)
		p0 = 1 - p1
		H = 0.0
		if p0 > 0: H -= p0 * math.log2(p0)
		if p1 > 0: H -= p1 * math.log2(p1)
		entropies.append((t, H))
		return entropies
		
		if __name__ == "__main__":
		history = run_rule_30(200)
		ent = sliding_entropy(history, 40)
		for t, H in ent:
		if t in [40, 80, 120, 158, 170, 200]:
		marker = " <- KRITISCHER PUNKT" if t == 158 else ""
		print(f"t = {t:3d}  H = {H:.5f}{marker}")
	\end{lstlisting}
	
	\subsubsection*{결과와 해석}
	
	스크립트를 실행하면 \(t = 158\)에서 엔트로피 \(H = 1.00000\)을 얻어, 중심 열이 임계 세대에서 정확히 완전한 유사 난수 생성기가 됨을 확인해 준다. 이 시점 이전에는 엔트로피가 점차 상승하며, 그 이후에는 포화된 채로 남는다. 이 거동은 해석적 표현 (\ref{eq:deltaS}) 및 YUCT 상수 \(\kappa_c = 1/3\)과 \(q = (3/2)^{1/3}\)로부터 유도된 \(t_{\mathrm{crit}}\) 값과 일치한다.
	
	이 실험은 매개변수 자유이다. 오직 YUCT의 기본 상수들과 규칙 30의 기본적 정의만을 사용한다. 예측된 포화점으로부터의 어떤 편차도 유효 오차 문턱 \(\kappa_c\)를 제약할 것이며, 따라서 보편적 오차 법칙의 직접적 검증을 제공할 것이다.
	
	\section{결론: 현실의 통일성}
	
	우리는 야쿠셰프 통일 협동 이론의 보편 상수(\(\beta = 2/3\), \(S_{\mathrm{odd}} = 1.2\))와 카오스 이론의 보편 상수(\(\alpha \approx 2.5029\), \(\delta \approx 4.6692\)) 사이의 형식적 대수적 연결을 제시하였다. 그 다리는 두 기둥 위에 놓여 있다.
	\begin{enumerate}
		\item 주기 배가의 재규격화군 이론으로부터의 엄밀한 해석적 관계 \(2\alpha^2 = \pi \delta\).
		\item YUCT에서 유도된 고정밀 근사 \(\pi \approx S_{\mathrm{odd}}\varphi^2\), 이는 원의 기하학을 협동의 기본 상수에 연결한다.
	\end{enumerate}
	
	결과 방정식 \(2\alpha^2 \approx (S_{\mathrm{odd}}\varphi^2)\delta\)는 창조(질서)의 상수와 파괴(카오스)의 상수가 독립적이 아님을 보여준다. 그것들은 황금비와 원에 의해 지배되는 단일의 통일된 수학적 현실의 두 측면이다.
	
	작지만 유한한 편차 \(\Delta\pi \approx 4.8\times10^{-5}\)는 수학적 결함이 아니라 \textbf{YUCT의 반증 가능한 예측}이다. 그것은 물리적 기하학의 입자성을 정량화하며 이론에 대한 직접적 실험적 접근을 제공한다. 부록 \(\Lambda\) 및 Ferra1.2와의 관련은 이 편차가 YUCT 틀의 견고한 특징이며 계층성 문제에서 협동 상 전이에 이르는 맥락에서 일관되게 나타남을 확인해 준다.
	
	이 발견은 심오한 함의를 가진다.
	\begin{itemize}
		\item \textbf{기초 물리학에 대해:} 파이겐바움 상수가 YUCT 상수와 마찬가지로 더 깊은 대수적 구조로부터 유도 가능할 수 있으며, \(\pi\)가 근본적이 아니라 창발적 상수임을 시사한다.
		\item \textbf{복잡계에 대해:} 고도로 협동된 계(뇌, 경제, 사회)가 언제 카오스적 불안정성의 문턱에 접근하는지 예측하기 위한 정량적 틀을 제공한다.
		\item \textbf{실험 과학에 대해:} 중력파 간섭계, 우주론적 서베이, 양자 간섭계를 사용한 구체적이고 가까운 미래의 검증을 제안하며, YUCT를 결정적으로 확인하거나 반박할 수 있을 것이다.
		\item \textbf{철학에 대해:} 질서와 카오스가 반대되는 것이 아니라 단일의 통일된 현실의 상보적 측면이라는 고래의 직관에 엄밀한 수학적 증명을 제공한다.
	\end{itemize}
	
	살아 있는 세포의 조화로운 협동에서 카오스적 흐름의 예측 불가능한 난류에 이르기까지, 현실의 통일성이 이제 하나의 우아한 방정식으로 표현된다. YUCT는 질서의 문법을 제공하고, 파이겐바움은 카오스의 문법을 제공하며, \(\pi\)와 \(\varphi\)는 그 둘이 쓰인 글자들이다. 편차 \(\Delta\pi\)는 이야기가 아직 끝나지 않았음을 알리는 문장 부호이다 – 그것은 우리가 현실의 협동적 직물에 대한 이해를 측정하고, 검증하고, 다듬도록 초대한다.
	
	\begin{thebibliography}{99}
		\bibitem{Yakushev2026} A.~V.~Yakushev, \emph{Yakushev Unified Coordination Theory (YUCT)}, Zenodo, DOI:10.5281/zenodo.18444599 (2026).
		\bibitem{AppendixAF} A.~V.~Yakushev, „Appendix AF: The Algebraic Framework of Reality in YUCT,“ in \emph{YUCT}, Zenodo (2026).
		\bibitem{AppendixFerra} A.~V.~Yakushev, „Appendix Ferra1.2: Universal Coordination Constants for Ordered and Disordered Phases,“ in \emph{YUCT}, Zenodo (2026).
		\bibitem{AppendixEhrenfest} A.~V.~Yakushev, „Appendix Ehrenfest: Coordination Interpretation of the Rotating Disk Paradox,“ in \emph{YUCT}, Zenodo (2026).
		\bibitem{AppendixLambda} A.~V.~Yakushev, „Appendix $\Lambda$: Resolution of the Hierarchy and Cosmological Constant Problems within YUCT,“ in \emph{YUCT}, Zenodo (2026).
		\bibitem{Feigenbaum1978} M.~J.~Feigenbaum, „Quantitative universality for a class of nonlinear transformations,“ \emph{Journal of Statistical Physics}, \textbf{19}(1), 25--52 (1978).
		\bibitem{Selvam2000} A.~M.~Selvam, „Universal characteristics of fractal fluctuations in prime number distribution,“ \emph{arXiv preprint}, physics/0008010 (2000).
		\bibitem{Briggs1992} J.~Briggs, \emph{Fractals: The Patterns of Chaos}. Simon \& Schuster (1992).
		\bibitem{Prigogine1977} I.~Prigogine and G.~Nicolis, \emph{Self-Organization in Nonequilibrium Systems}. Wiley (1977).
		\bibitem{Prigogine1984} I.~Prigogine and I.~Stengers, \emph{Order out of Chaos}. Bantam (1984).
		\bibitem{Rayleigh1916} Lord Rayleigh, „On convection currents in a horizontal layer of fluid,“ \emph{Philosophical Magazine}, \textbf{32}, 529--546 (1916).
		\bibitem{Chandrasekhar1961} S.~Chandrasekhar, \emph{Hydrodynamic and Hydromagnetic Stability}. Oxford (1961).
		\bibitem{Wolfram2002} S.~Wolfram, \emph{A New Kind of Science}, Wolfram Media (2002).
	\end{thebibliography}
	
\end{document}